%% =====================================================================
%%  paper2_introduction_v2.tex
%%
%%  Revised Introduction (Section 1) for
%%    "Delay-Integrated Cascading Failure Analysis in Cyber-Physical
%%     Power Systems: A Closed-Loop Feedback Perspective"
%%
%%  Drop-in replacement for the \section{Introduction} block currently
%%  located at lines 120--218 of paper2_delay_cascade.tex.  Compiles
%%  against the Elsevier cas-dc class already used by paper2.
%%
%%  Design rationale (so future revisions stay consistent):
%%    * Frames the contribution against THREE concrete tiers of prior
%%      work, not a vague "most papers ignore delay" statement, so a
%%      reviewer cannot dismiss the gap as already closed.
%%    * Every tier names representative classic + recent (2022--2024)
%%      references, addressing the user's request to include both
%%      seminal and up-to-date sources.
%%    * Each contribution bullet is anchored to a specific industrial
%%      standard (NERC PRC-002-2, NASPI P95, IEC 61850 PRP,
%%      IEEE C37.118.2, NERC CIP-012, NERC PRC-006-5) so the model
%%      parameters are externally defensible, not arbitrary.
%%    * Preserves the closed-loop feedback narrative, the
%%      \(R_{1}\)/\(R_{3}\) metric pair, and the \eta^{+} mechanism
%%      decomposition that the rest of the paper relies on.
%%
%%  New \citep keys introduced here are defined in
%%    paper2_refs_v2_additions.bib
%%  and should be \input or \bibinput-merged into paper2_refs.bib
%%  before final compilation.
%% =====================================================================

\section{Introduction}\label{sec:intro}

The deep integration of physical infrastructure with high-speed
communication networks has transformed traditional power grids into
\emph{cyber-physical power systems} (CPPS), in which a physical power
layer and a cyber communication layer operate as a tightly coupled whole
\citep{Rinaldi2001,Ouyang2014,Kundur2017CPPS}.  Under nominal conditions
this coupling improves dispatch flexibility, situational awareness and
wide-area control performance.  Under stress, however, it opens a
bidirectional failure channel: a disturbance originating in either layer
can propagate to the other and trigger cascading failures of a magnitude
neither layer would experience in isolation
\citep{Buldyrev2010,Gao2012,Guo2017AppliedEnergy}.

\subsection*{From a binary cyber layer to a delay-quality cyber layer}

A rich body of work has investigated cascading failures through
topological models, flow-redistribution rules and attack-vulnerability
analyses
\citep{Motter2002,Dobson2007,Brummitt2012,Schneider2011}.
These studies established \emph{how} structural interdependence amplifies
failure propagation, but they treat the communication layer as a binary
entity---a link is either functional or failed---without modelling the
\emph{quality of service} that surviving links provide.  In particular,
end-to-end latency, defined as the time elapsed between issuing a
measurement or a control command and its arrival at the destination, is
either absent from the cascade dynamics or appears only as a post-hoc
observation (``delay increased after the failures'') rather than as a
causal variable that degrades generator control \emph{during} the
cascade itself \citep{Milano2012}.

This omission is substantive, not cosmetic.  Industrial reliability
standards explicitly bound end-to-end latency for protection-grade
control: NERC PRC-002-2 requires fault recording and event reporting
within sub-second windows \citep{NERC_PRC002}, NASPI synchrophasor
deployment guidelines report a 95th-percentile end-to-end latency in the
$200$--$500\,\textrm{ms}$ band for utility WAMS
\citep{NASPI_LatencyGuidelines}, and IEEE~C37.118.2 fixes the report
formatting and time-alignment budget for PMU streams
\citep{IEEE_C37_118_2}.  Once any one of these budgets is consumed by
queueing in a stressed cyber topology, the supervisory control loop
operates on stale state and lagged actuation, and the closed-loop
controllability of generation declines well before any second physical
asset has tripped.

\subsection*{Three tiers of prior work, and a missing fourth}

To position the present paper precisely, we group existing CPPS
cascading-failure studies that touch on communication latency into three
tiers, summarised in Table~\ref{tab:prior_tiers}, and identify the tier
that has remained empty.

\begin{table*}[!t]
\centering
\caption{Three tiers of prior CPPS cascading-failure studies w.r.t.\ the
treatment of communication delay, and the gap addressed by this work.}
\label{tab:prior_tiers}
\small
\begin{tabular}{@{}p{0.16\linewidth}p{0.36\linewidth}p{0.40\linewidth}@{}}
\toprule
\textbf{Tier} & \textbf{How delay is modelled} & \textbf{Representative work} \\
\midrule
\textbf{T0 — Instantaneous cyber} &
Cascade is driven by topology and DC/AC overload only;
``commands propagate without delay'' is an explicit or implicit
assumption. &
\citet{Motter2002,Dobson2007,Brummitt2012,Buldyrev2010,Gao2012} \\[2pt]
\textbf{T1 — Single lumped delay~$\tau$} &
The whole cyber chain is collapsed into a constant scalar delay or a
constant message-loss probability, applied uniformly to every
control action. &
\citet{Cai2016TSG,Falahati2014TSG,Guo2017AppliedEnergy,Huang2015} \\[2pt]
\textbf{T2 — One-channel variable delay} &
Delay is allowed to vary, but only along a \emph{single} axis: queueing
on transmission links, link overload, or PMU jitter for state
estimation. The other channels of the cyber stack remain idealised. &
\citet{Xin2024CPPSframework,LinkOverload2024Electronics,Wang2022CPPS,LaiCPPS2019} \\[2pt]
\textbf{T3 — \emph{This work}: full-stack, standard-anchored,
closed-loop} &
End-to-end delay is decomposed into measurement, edge processing,
transmission, and actuation/UFLS components; each component is bounded
by a named industrial standard; the resulting heterogeneous $\boldsymbol{\tau}$
is fed through a triple-mechanism efficacy function
$\eta^{+}=\Phi_{\mathrm{loss}}\!\cdot\!\Phi_{\mathrm{sat}}\!\cdot\!
\Phi_{\mathrm{crit}}$ inside every cascade round, closing the
delay~$\to$~control~$\to$~overload~$\to$~topology~$\to$~delay loop. &
--- (this paper) \\
\bottomrule
\end{tabular}
\end{table*}

Tier~T0 omits the cyber service quality entirely; tier~T1 makes
the delay penalty independent of \emph{where} in the cyber stack the
bottleneck actually sits, so every mitigation looks equally effective;
tier~T2 captures one channel realistically but misses the interaction
between channels.  None of the three tiers offers a model in which
(i)~latency originates from the same physical mechanisms a NASPI
audit would record, (ii)~the latency feeds back into the
\emph{next} round of the cascade rather than being applied as a
post-hoc rescaling, and (iii)~the resulting damage is observable on
both an aggregate load metric and a generator-fidelity metric.  Closing
this gap is the objective of the present paper.

\subsection*{The closed-loop mechanism we embed}

In a closed-loop control architecture, stale measurements cause the
control center to act on outdated states while lagged execution causes
corrective actions to arrive after the system has drifted further.  Both
effects reduce the effective generation output relative to its scheduled
reference.  When this reduction is embedded \emph{before} the power-flow
redistribution that determines which branches overload, a
self-reinforcing feedback loop emerges:
\begin{quote}
\emph{Delay reduces effective generation
$\;\to\;$ altered power flows cause additional overloads
$\;\to\;$ overloaded branches trip
$\;\to\;$ structural failures degrade the cyber topology
$\;\to\;$ the degraded topology lengthens paths and amplifies delay
$\;\to\;$ further generation reduction in the next round
$\;\to\;\cdots$.}
\end{quote}
This loop converts communication latency from a passive symptom of
distress into an active driver of further damage.  It is by
construction invisible to T0/T1 models, and only partially visible to
T2 models, because in those models all delay scenarios share a
common structural cascade and delay merely rescales final performance
metrics \citep{LaiCPPS2019}.  Our prior work
\citep{ChenTu2026} established the underlying CPPS framework with
distributed control centers, a Functional
Viability-Constrained Island Detection (FVC-ID) mechanism and a
Cohesion-Propagation Importance Metric (CPIM) for control-center
placement, but did not model the influence of communication delay on
generator control during the cascade.  The present paper fills that
gap.

\subsection*{Contributions}

The contributions of this work are as follows.

\begin{enumerate}
\item \textbf{A delay-integrated cascading failure framework with
  closed-loop feedback.}  Delay computation is embedded inside every
  cascade round and scales generator output \emph{before} the
  DC power-flow redistribution.  Each delay scenario therefore
  produces a genuinely different cascade trajectory---different
  overload patterns, different failure sequences, different terminal
  states---rather than a rescaled copy of a common trajectory, in
  contrast to T0--T2 models above.

\item \textbf{A standard-anchored, role-aware, four-component delay
  decomposition.}  End-to-end latency is decomposed into measurement
  ($\tau_m$), edge processing ($\tau_e$), transmission ($\tau_t$) and
  actuation/UFLS ($\tau_a$) components, each parameterised against a
  named industrial reference: NASPI synchrophasor latency
  guidelines~\citep{NASPI_LatencyGuidelines} and
  IEEE~C37.118.2~\citep{IEEE_C37_118_2} for $\tau_m$ and~$\tau_t$,
  IEC~61850 PRP/HSR redundancy~\citep{IEC_61850_PRP} for the
  $k_{\max}=3$ redundancy stack used inside~$\Phi_{\mathrm{loss}}$,
  NERC~CIP-012~\citep{NERC_CIP012} for control-center traffic
  protection, and NERC PRC-002-2~\citep{NERC_PRC002} for the upper
  envelope used to define the \texttt{heavy} scenario
  ($\tau_{\mathrm{tot}}\!\approx\!594\,\mathrm{ms}\!\approx\!0.6\times$
  PRC-002-2 budget).  Uplink (measurement) and downlink (command)
  paths are computed separately, with role-dependent service times
  for control centers versus relay nodes, capturing the asymmetric
  processing behaviour inherent in SCADA/EMS architectures.

\item \textbf{A three-mechanism control-efficacy model
  $\eta^{+}=\Phi_{\mathrm{loss}}\!\cdot\!\Phi_{\mathrm{sat}}\!\cdot\!
  \Phi_{\mathrm{crit}}$.}  Delay-induced robustness loss is
  decomposed into (i)~redundancy-attenuated message loss
  ($\Phi_{\mathrm{loss}}$, with $k_{\max}$-fold IEC~61850~PRP backup),
  (ii)~controller saturation under elevated $\tau_m+\tau_e$
  ($\Phi_{\mathrm{sat}}$), and (iii)~normalised critical collapse
  near the per-machine $\tau_{\mathrm{crit},i}$
  ($\Phi_{\mathrm{crit}}$, normalised so that
  $\eta^{+}(\tau\!=\!0)\!=\!1$ exactly).  The three mechanisms are
  orthogonal in parameter space, allowing a reviewer to attribute any
  observed delay penalty to a specific physical channel.

\item \textbf{Coupled UFLS over-shedding under delay.}  In contrast
  to the standard NERC~PRC-006-5 single-stage UFLS rule
  \citep{NERC_PRC006}, we model that elevated $\tau_m+\tau_e$
  forces the load-shedding scheme to over-shed by a delay-dependent
  factor $\gamma(\tau)$ before stabilisation, capped at a
  configurable $\overline{\phi}_{\max}$.  This is, to our knowledge,
  the first cascading-failure model in which UFLS depth is an
  endogenous function of cyber-layer latency rather than a fixed
  rule.

\item \textbf{Twin resilience metrics that reveal hidden degradation.}
  The cascade is evaluated jointly on $R_1$, a supply-conservation
  ratio, and $R_3$, a capacity-weighted dispatch-fidelity NRMSE.  We
  show that delay damage is asymmetrically concentrated at
  moderate-to-high tolerance margins---precisely where a delay-free
  analysis predicts safety---and that $R_3$ exposes hidden functional
  degradation that aggregate load-retention metrics miss.  A
  delay-penalty heatmap on the (tolerance, cascade-round) plane
  identifies a narrow critical intervention window in the early
  rounds.
\end{enumerate}

\noindent
The combined effect is that delay-free and single-$\tau$ robustness
models are not merely imprecise; they are systematically
\emph{misleading}, providing the strongest reassurance in the operating
regimes that harbour the greatest latent vulnerability.

\subsection*{Paper organisation}

The remainder of this paper is organised as follows.
Section~\ref{sec:model} presents the CPPS model, including the
four-component communication delay formulation and the
betweenness-assortative interdependent coupling.
Section~\ref{sec:cascade} describes the delay-integrated cascading
failure dynamics, the $\eta^{+}$ control-efficacy function, the
delay-coupled UFLS rule, and the five-level delay severity scenarios.
Section~\ref{sec:experiment} reports Monte~Carlo experiments on the
IEEE~39-bus test system coupled to a Barab\'asi--Albert scale-free
communication network, covering 1\,000 coupling realisations, eleven
tolerance levels and five delay scenarios, and reports the asymmetric
delay-penalty regime, the heavy-tailed collapse distributions, the
hidden $R_3$ degradation, the round-level intervention window, and a
sensitivity analysis of mitigation actions.
Section~\ref{sec:conclusion} concludes the paper.
